\begin{khmer}
    
ការរាប់គ្រឿងភ្ជាប់លោហៈខ្នាតតូចប្រកបដោយភាពត្រឹមត្រូវ និងក្នុងពេលវេលាជាក់ស្តែងនៅលើខ្សែក្រវាត់ចលនាឧស្សាហកម្មល្បឿនលឿន នៅតែជាបញ្ហាប្រឈមដ៏សំខាន់មួយក្នុងវិស័យភស្តុភារផលិតកម្ម។ ការត្រួតពិនិត្យដោយភ្នែកមនុស្សតម្រូវឱ្យមានការផ្តោតអារម្មណ៍ខ្ពស់ ងាយកើតមានកំហុស និងមិនមានប្រសិទ្ធភាពផ្នែកសេដ្ឋកិច្ច ខណៈដែលការប៉ាន់ស្មានតាមទម្ងន់បែបប្រពៃណីមិនអាចផ្តល់នូវកម្រិតលម្អិតដែលចាំបាច់សម្រាប់ការតាមដានខ្សែសង្វាក់ផ្គត់ផ្គង់ទំនើបបាន។ លើសពីនេះ ដំណោះស្រាយស្វ័យប្រវត្តិកម្មដែលមានស្រាប់នៅតែប្រឈមនឹងបញ្ហាដូចជា ពន្លឺចាំងផ្លាតលើផ្ទៃលោហៈដែលមានលក្ខណៈចាំង ការត្រួតស៊ីគ្នាយ៉ាងកកកុញរបស់ផ្នែក និងតម្លៃគណនាខ្ពស់របស់ម៉ូដែលចាប់សញ្ញាដែលមានភាពត្រឹមត្រូវខ្ពស់ ដែលបញ្ហាទាំងនេះរួមគ្នារារាំងដល់ការដាក់ពង្រាយប្រព័ន្ធដែលអាចទុកចិត្តបាន និងដំណើរការជាបន្តបន្ទាប់នៅកម្រិត \textenglish{edge} ក្នុងរោងចក្រ។ \\

និក្ខេបបទនេះបង្ហាញនូវក្របខ័ណ្ឌ \textenglish{Edge-AI} ពេញលេញ និងមិនពឹងផ្អែកលើពពក សម្រាប់ការចាប់សញ្ញា ការចាត់ថ្នាក់ និងការរាប់គ្រឿងភ្ជាប់លោហៈនៅលើខ្សែក្រវាត់ចលនាឧស្សាហកម្មក្នុងពេលវេលាជាក់ស្តែង។ ស្នូលនៃប្រព័ន្ធគឺ \textenglish{YOLO11-GED} ដែលជាស្ថាបត្យកម្មចាប់សញ្ញាវត្ថុទម្ងន់ស្រាល ដែលត្រូវបានពង្រឹងដោយម៉ូឌុលថ្មី \textenglish{EMAc2f}។ ម៉ូឌុលនេះរួមបញ្ចូលយន្តការបំពេញបន្ថែមគ្នាចំនួនបី៖ \textenglish{Efficient Multi-Attention} សម្រាប់កាត់បន្ថយឥទ្ធិពលនៃភាពមិនប្រក្រតីដែលបណ្តាលមកពីពន្លឺចាំងផ្លាត, ការប្រើប្រាស់លក្ខណៈពិសេសឡើងវិញផ្អែកលើ \textenglish{GhostNet} សម្រាប់ការតំណាងវត្ថុខ្នាតតូចប្រកបដោយប្រសិទ្ធភាពនៃប៉ារ៉ាម៉ែត្រ និង \textenglish{Dynamic Upsampling} សម្រាប់រក្សាភាពត្រឹមត្រូវនៃព័ត៌មានលំហក្នុងតម្លៃគណនាទាប។ ដើម្បីបង្កើតមូលដ្ឋានប្រៀបធៀបប្រកបដោយភាពម៉ឺងម៉ាត់ ការកំណត់រចនាសម្ព័ន្ធម៉ូដែលចាប់សញ្ញាចំនួនប្រាំបួន ដែលគ្របដណ្តប់លើជំនាន់ \textenglish{YOLO} ចំនួនបី (\textenglish{YOLOv8n/s/m, YOLOv11n/s/m, YOLOv26n/s/m}) ត្រូវបានធ្វើតេស្ត និងវាស់ស្ទង់លើសំណុំទិន្នន័យដែលបង្កើតឡើងជាពិសេសសម្រាប់រូបភាពគ្រឿងភ្ជាប់ ដែលបានថតក្រោមលក្ខខណ្ឌខ្សែក្រវាត់ចលនាជាក់ស្តែង រួមមានពន្លឺប្រែប្រួល ភាពព្រិលដោយសារចលនា និងការប្រមូលផ្តុំយ៉ាងកកកុញនៃគ្រឿងភ្ជាប់។ \\

ម៉ូដែលដែលបានបណ្តុះបណ្តាលត្រូវបានដាក់ពង្រាយលើវេទិកាកុំព្យូទ័របង្កប់ \textenglish{NVIDIA Jetson} ជាមួយនឹងការបង្កើនល្បឿនដោយ \textenglish{TensorRT} និងភ្ជាប់ជាមួយកាមេរ៉ាជម្រៅ \textenglish{Intel RealSense}។ ប្រព័ន្ធនេះសម្រេចបានគោលដៅនៃ \textenglish{mAP@0.5 ≥ 0.95} នៅល្បឿន \textenglish{≥ 30 FPS} ដោយមិនពឹងផ្អែកលើសេវាកម្មពពកឡើយ។ ជុំវិញម៉ាស៊ីនចក្ខុវិស័យ ប្រព័ន្ធផលិតកម្មពេញលេញមួយត្រូវបានអនុវត្ត ដែលរួមមានសេវាកម្មសន្និដ្ឋានម៉ូដែល \textenglish{Python}, \textenglish{Express.js REST API}, មូលដ្ឋានទិន្នន័យអចិន្ត្រៃយ៍សម្រាប់រក្សាទុកកំណត់ត្រាសវនកម្ម ស្រទាប់ \textenglish{Redis} សម្រាប់ទិន្នន័យតេឡេមេទ្រីក្នុងពេលវេលាជាក់ស្តែង និងចំណុចប្រទាក់សម្រាប់ប្រតិបត្តិករចំនួនពីរ រួមមានផ្ទាំងគ្រប់គ្រងតាមបណ្តាញ (\textenglish{Web Dashboard}) និងកម្មវិធីកុំព្យូទ័រលើតុ (\textenglish{Desktop Application})។ ប្រព័ន្ធដែលបានរួមបញ្ចូលនេះត្រូវបានផ្ទៀងផ្ទាត់លើឧបករណ៍សាកល្បងខ្សែក្រវាត់ចលនាជាក់ស្តែង ដោយបង្ហាញពីសមត្ថភាពរាប់ដែលមានស្ថេរភាព និងពន្យារពេលទាប ព្រមទាំងកាត់បន្ថយតម្រូវការក្នុងការរាប់ឡើងវិញដោយមនុស្ស និងកំហុសដែលបណ្តាលមកពីមនុស្ស។ \\

ក្រៅពីលទ្ធផលដែលអាចអនុវត្តបានក្នុងបរិបទឧស្សាហកម្ម ការងារនេះក៏បានចូលរួមចំណែកនូវវិធីសាស្ត្រដែលអាចធ្វើឡើងវិញបានចាប់ពីការរចនារហូតដល់ការដាក់ពង្រាយ ព្រមទាំងការសិក្សា \textenglish{ablation} ជាបរិមាណ ដែលបំបែកឥទ្ធិពលរបស់សមាសធាតុនីមួយៗនៃ \textenglish{EMAc2f}។ លើសពីនេះ វាបានផ្តល់នូវភស្តុតាងដែលបានកត់ត្រាថា ការកែសម្រួលស្ថាបត្យកម្មឱ្យសមស្របតាមបញ្ហាអាចបង្កើនសមត្ថភាពរំលឹកវត្ថុខ្នាតតូច (\textenglish{small-object recall}) ក្រោមកម្រិតធនធានគណនាដ៏តឹងរ៉ឹងរបស់ប្រព័ន្ធ \textenglish{edge} ដោយមិនប៉ះពាល់ដល់ល្បឿនដំណើរការរបស់ប្រព័ន្ធ។
\end{khmer}
